\chapter{Curve Sketching: The Complete Discussion}

Everything needed to draw a curve has now been assembled in pieces:
domains and asymptotes in the functions volume, limits in the analysis
volume, the first derivative for monotonicity, the second for curvature.
What is missing is the habit of running all of it in one fixed order and
finishing with a picture. That is the \de{Kurvendiskussion}, and it is worth
drilling as a single routine for two reasons.

First, examinations ask for it directly. Second, and more usefully, the
routine is a self-checking mechanism. A sign table that says the function
falls on $(0,2)$ and a computed local maximum at $x=2$ cannot both be right;
the contradiction surfaces before the picture is drawn. A curve sketch is
the cheapest way there is to catch an algebra slip.

The chapter runs the protocol on rational, exponential, logarithmic and root
functions, and closes with the discussion an economist actually performs:
locating the inflection of a total cost curve, the minimum of average cost,
and the output that maximises profit --- all of them the same routine
applied to a curve with a name.

\section{Core Rules}

\begin{rulebox}{The protocol}
Always in this order, because each step feeds the next:
\begin{enumerate}[leftmargin=1.8em,itemsep=0.35ex]
  \item \textbf{Domain.} Denominators, even roots, logarithms. Note which
        points are excluded --- they are the candidates for poles.
  \item \textbf{Symmetry.} $f(-x)=f(x)$ even, $f(-x)=-f(x)$ odd. Either one
        halves the work.
  \item \textbf{Intercepts and roots.} $f(0)$, and the solutions of
        $f(x)=0$ with their multiplicities.
  \item \textbf{Behaviour at the edges.} Limits at $\pm\infty$ and at every
        excluded point, one side at a time. This produces the
        \textbf{asymptotes}: vertical at poles, horizontal or oblique at
        infinity.
  \item \textbf{First derivative.} Solve $f'(x)=0$ for the critical points,
        build a \textbf{sign table} for $f'$, and read off intervals of
        increase and decrease and the type of each extremum.
  \item \textbf{Second derivative.} Solve $f''(x)=0$, build a sign table
        for $f''$, and read off convex and concave intervals and the
        \textbf{inflection points}.
  \item \textbf{Value table} for any remaining doubt, then \textbf{draw}:
        asymptotes first as dashed guides, then the special points, then
        join them respecting monotonicity and curvature.
\end{enumerate}
\end{rulebox}

\begin{rulebox}{Sign tables}
A sign table is the bookkeeping device that prevents guesswork. List every
point where the expression can change sign --- its roots and its poles ---
in increasing order; they cut the domain into intervals. On each interval
the sign is constant, so \emph{one} test value settles it. For
$f'(x)=3x(x-2)$:
\[
  \begin{array}{c|ccccc}
    x & (-\infty,0) & 0 & (0,2) & 2 & (2,\infty)\\\hline
    f' & + & 0 & - & 0 & +\\
    f & \nearrow & \text{max} & \searrow & \text{min} & \nearrow
  \end{array}
\]
A sign change $+\to-$ is a local maximum, $-\to+$ a local minimum, and
\emph{no} change means no extremum, however much the derivative vanishes
(as at $x=0$ for $f=x^{3}$).
\end{rulebox}

\begin{rulebox}{First- and second-derivative tests}
Let $f'(x_{0})=0$.
\begin{itemize}[leftmargin=1.6em,itemsep=0.3ex]
  \item \textbf{First-derivative test:} the sign change of $f'$ across
        $x_{0}$ decides the type. Always applicable, and the only reliable
        test when $f''(x_{0})=0$.
  \item \textbf{Second-derivative test:} $f''(x_{0})>0$ gives a local
        minimum, $f''(x_{0})<0$ a local maximum. If $f''(x_{0})=0$ the test
        is \textbf{silent} --- fall back on the sign table. Compare
        $x^{4}$ (minimum), $-x^{4}$ (maximum) and $x^{3}$ (neither), all
        three with $f'=f''=0$ at the origin.
  \item \textbf{Inflection} at $x_{0}$ requires $f''$ to \emph{change
        sign} there. $f''(x_{0})=0$ alone is not enough: $f=x^{4}$ has
        $f''(0)=0$ and no inflection.
\end{itemize}
\end{rulebox}

\begin{rulebox}{Asymptotes, recalled}
For $f=p/q$ with $p,q$ polynomials, after cancelling common factors:
vertical asymptote at each zero of $q$; horizontal $y=0$ if
$\deg p<\deg q$; horizontal $y=$ ratio of leading coefficients if
$\deg p=\deg q$; and if $\deg p=\deg q+1$, polynomial division gives
\[
  f(x)=\underbrace{mx+c}_{\text{oblique asymptote}}+\frac{\text{remainder}}{q(x)} .
\]
For non-rational functions the same information comes from limits directly:
$\lim_{x\to\pm\infty}f(x)$, and for a possible oblique asymptote
$m=\lim_{x\to\infty}f(x)/x$ followed by $c=\lim_{x\to\infty}(f(x)-mx)$.
\end{rulebox}

\begin{gallery}
\hspace*{\fill}%
\begin{tikzpicture}
  \begin{axis}[fnaxis, width=8.4cm, height=7.4cm,
      xmin=-4.4, xmax=5.4, ymin=-6.4, ymax=10.4,
      xtick={-4,-2,2,4}, ytick={-4,4,8}]
    \addplot[gray!55, dashed, thin, samples=2, forget plot]
      coordinates {(1,-6.2)(1,10.2)};
    \addplot[gray!55, dashed, thin, samples=2, domain=-4.2:5.2, forget plot]{x+1};
    \addplot[domain=-4.2:0.86, samples=201]{x^2/(x-1)};
    \addplot[domain=1.14:5.2, samples=201]{x^2/(x-1)};
    \addplot[only marks, mark=*, mark size=1.3pt, black, forget plot]
      coordinates {(0,0) (2,4)};
    \node[font=\tiny, anchor=north east] at (axis cs:-0.1,-0.15) {$(0,0)$};
    \node[font=\tiny, anchor=west] at (axis cs:2.15,4) {$(2,4)$};
  \end{axis}
\end{tikzpicture}%
\hspace*{\fill}
\gcap{The finished discussion of $f(x)=x^{2}/(x-1)$: pole $x=1$, oblique
asymptote $y=x+1$, local maximum $(0,0)$ below the local minimum $(2,4)$.}
\end{gallery}

\begin{rulebox}{The economist's curve discussion}
Three facts fall straight out of the protocol applied to a total cost
function $C(q)$, and they are worth knowing as results:
\begin{itemize}[leftmargin=1.6em,itemsep=0.4ex]
  \item Marginal cost is $MC=C'$. The \textbf{inflection point of $C$} is
        exactly where $MC$ is smallest, since $MC'=C''$ changes sign there.
        The S-shaped cost curve turns from concave to convex at its most
        efficient scale.
  \item Average cost is $AC=C(q)/q$. Differentiating,
        \[
          AC'=\frac{C'(q)q-C(q)}{q^{2}}=\frac{MC-AC}{q},
        \]
        so $AC'=0$ exactly when $MC=AC$: \textbf{marginal cost cuts average
        cost at the minimum of average cost}. That is a theorem, not a
        coincidence of a particular example.
  \item Profit at price $p$ is $\pi=pq-C(q)$, so $\pi'=p-MC$ and the
        optimum has $p=MC$, confirmed as a maximum by $\pi''=-C''<0$,
        i.e.\ by cost being convex there.
\end{itemize}
\end{rulebox}

\section{Worked Example}

\begin{worked}{The complete discussion of $f(x)=\dfrac{x^{2}}{x-1}$}
\wstep{Step 1. Domain, and the reason to find it first.}
The denominator vanishes at $x=1$, so
$D_{f}=\R\setminus\{1\}$. Every later step --- sign tables, monotonicity,
convexity --- must be stated \emph{per interval}, and $x=1$ is the wall
between them.

\wstep{Step 2. Symmetry and roots.}
$f(-x)=\frac{x^{2}}{-x-1}$, which is neither $f(x)$ nor $-f(x)$: no
symmetry. The numerator vanishes at $x=0$ with multiplicity two, so the
curve \textbf{touches} the axis at the origin rather than crossing it, and
$f(0)=0$.

\wstep{Step 3. Behaviour at the pole, from both sides.}
Near $x=1$ the numerator is close to $1>0$, so the sign of $f$ is the sign
of $x-1$:
\[
  x\to1^{-}:\ \frac{1^{+}}{0^{-}}\to-\infty,
  \qquad
  x\to1^{+}:\ \frac{1^{+}}{0^{+}}\to+\infty .
\]
\wtrap{Take one-sided limits at every pole. Writing ``$f\to\infty$ at
$x=1$'' loses the sign information and makes the sketch impossible to draw
correctly.}

\wstep{Step 4. Behaviour at infinity: divide, do not guess.}
Numerator degree $2$ exceeds denominator degree $1$ by exactly one, so there
is an \textbf{oblique} asymptote. Polynomial division gives it:
\[
  f(x)=x+1+\frac{1}{x-1} .
\]
So $y=x+1$ is the asymptote, and the remainder $\frac{1}{x-1}$ says how the
curve sits relative to it: \emph{above} for $x>1$, \emph{below} for $x<1$.
\wtrap{Horizontal asymptote when the degrees are equal or the numerator's is
smaller; oblique when the numerator's degree is exactly one larger; neither
when it is larger by two or more. And the division is the whole method ---
guessing $y=x$ from the leading behaviour misses the $+1$.}

\wstep{Step 5. First derivative and critical points.}
Differentiating the divided form is easiest:
\[
  f'(x)=1-\frac{1}{(x-1)^{2}}=\frac{(x-1)^{2}-1}{(x-1)^{2}}=\frac{x(x-2)}{(x-1)^{2}} .
\]
Zeros at $x=0$ and $x=2$; the denominator is positive wherever it is
defined, so the sign of $f'$ is the sign of $x(x-2)$.
\begin{center}\footnotesize
\begin{tabular}{lcccc}
  \hline
  interval & $x<0$ & $0<x<1$ & $1<x<2$ & $x>2$\\
  \hline
  $f'$ & $+$ & $-$ & $-$ & $+$\\
  $f$ & rising & falling & falling & rising\\
  \hline
\end{tabular}
\end{center}

\wstep{Step 6. Classify, and notice something strange.}
$f''(x)=\dfrac{2}{(x-1)^{3}}$, so $f''(0)=-2<0$ and $f''(2)=2>0$:
\[
  \text{local maximum at }(0,0),
  \qquad
  \text{local minimum at }(2,4).
\]
The local maximum lies \textbf{below} the local minimum. That is impossible
for a continuous function --- and possible here only because the pole at
$x=1$ separates them. The two branches never see each other.
\wtrap{Local extrema are local. A local maximum of $0$ coexisting with a
local minimum of $4$ is not an arithmetic error, and neither point is a
global extremum: $f$ is unbounded in both directions.}

\wstep{Step 7. Convexity, and the absent inflection.}
$f''=\frac{2}{(x-1)^{3}}$ is never zero, so there is \textbf{no inflection
point}. It is negative for $x<1$ and positive for $x>1$:
concave on $(-\infty,1)$, convex on $(1,\infty)$. The change of curvature
happens at the pole, where the function does not exist --- so it is not an
inflection.
\wtrap{A change of convexity across a point of the domain is an inflection;
across a gap in the domain it is nothing at all. Check that the candidate
point is in the domain before naming it.}

\wstep{Step 8. Draw it.}
\begin{center}
\begin{tikzpicture}
\begin{axis}[fnaxis, width=8.4cm, height=6.0cm,
             xmin=-4.2, xmax=6.2, ymin=-9, ymax=13,
             xtick={-3,-1,1,2,4}, ytick={4,8},
             restrict y to domain=-40:40]
  \addplot[domain=-4:0.88, samples=200] {x^2/(x-1)};
  \addplot[domain=1.12:6, samples=200] {x^2/(x-1)};
  \addplot[accent!40, thick, dashed, samples=2, domain=-4:6, forget plot]
    {x+1};
  \addplot[gray!55, thin, dashed, samples=2, forget plot]
    coordinates {(1,-9) (1,13)};
  \addplot[only marks, mark=*, mark size=1.2pt, accent!70, forget plot]
    coordinates {(0,0) (2,4)};
\end{axis}
\end{tikzpicture}
\end{center}
Everything checks against the figure: the touch at the origin, the two
branches, the pole at $x=1$, the asymptote $y=x+1$ approached from below on
the left and from above on the right, and the maximum sitting lower than the
minimum.

\wstep{Step 9. The protocol, in the order that saves work.}
\begin{center}\footnotesize
\begin{tabular}{ll}
  \hline
  1. domain & 6. sign table for $f'$\\
  2. symmetry & 7. extrema, classified by $f''$\\
  3. roots and intercepts & 8. $f''$, convexity, inflections\\
  4. poles, one-sided limits & 9. sketch\\
  5. asymptotes (divide!) & 10. check the sketch against 1--8\\
  \hline
\end{tabular}
\end{center}
Step~10 is not optional. A sketch that contradicts the sign table means one
of them is wrong, and finding out which takes a minute --- far less than
losing the marks.
\end{worked}

\section{Practice Problems}

In problems 1--8 give the intervals of increase and decrease and all local
extrema, using a sign table.

\begin{problems}
  \item $f(x)=x^{3}-12x$
  \item $f(x)=x^{4}-8x^{2}$
  \item $f(x)=x\mathrm{e}^{-x}$
  \item $f(x)=\dfrac{x}{x^{2}+1}$
  \item $f(x)=x-\ln x$ on $x>0$
  \item $f(x)=x^{3}$ --- and explain why $f'(0)=0$ yields no extremum.
  \item $f(x)=\dfrac{x^{2}}{x-1}$
  \item $f(x)=\sqrt{x}-x$ on $x\ge0$
\end{problems}

In problems 9--14 give the convex and concave intervals and all inflection
points.

\begin{problems}[resume]
  \item $f(x)=x^{3}-3x^{2}$
  \item $f(x)=x^{4}-4x^{2}$
  \item $f(x)=\mathrm{e}^{-x^{2}}$
  \item $f(x)=x^{4}$ --- and explain why $f''(0)=0$ gives no inflection.
  \item $f(x)=\dfrac{\ln x}{x}$ on $x>0$
  \item $f(x)=x^{2}\mathrm{e}^{-x}$
\end{problems}

In problems 15--20 give every asymptote.

\begin{problems}[resume]
  \item $f(x)=\dfrac{x-1}{x+2}$
  \item $f(x)=\dfrac{1}{x^{2}-4}$
  \item $f(x)=\dfrac{x^{2}}{x-1}$
  \item $f(x)=\dfrac{x^{3}}{x^{2}+1}$
  \item $f(x)=\dfrac{\ln x}{x}$
  \item $f(x)=x+\mathrm{e}^{-x}$
\end{problems}

Problems 21--32 ask for a \emph{complete} discussion: domain, symmetry,
intercepts, asymptotes and edge behaviour, monotonicity and extrema,
curvature and inflections, and a sketch.

\begin{problems}[resume]
  \item $f(x)=x^{3}-3x^{2}$
  \item $f(x)=x^{4}-4x^{2}$
  \item $f(x)=\dfrac{x}{x^{2}+1}$
  \item $f(x)=\dfrac{x^{2}}{x-1}$
  \item $f(x)=\mathrm{e}^{-x^{2}}$
  \item $f(x)=x\mathrm{e}^{-x}$
  \item $f(x)=\dfrac{\ln x}{x}$
  \item $f(x)=\dfrac{x-1}{x+2}$
  \item $f(x)=x^{2}\mathrm{e}^{-x}$
  \item $f(x)=\dfrac{1}{x^{2}-4}$
  \item $f(x)=\sqrt{x}-x$
  \item $f(x)=x^{3}-3x+2$
\end{problems}

\begin{problems}[resume]
  \item A function satisfies $f'>0$ and $f''<0$ throughout an interval.
        Describe the shape of its graph there in one sentence.
  \item The graph of $f'$ is a parabola opening upwards with roots at
        $x=-1$ and $x=3$. What does $f$ do, and where are its extrema?
  \item The graph of $f'$ is a straight line through the origin with
        positive slope. What is $f$, up to a constant, and what is its
        shape?
  \item $f$ is differentiable with $f(1)=2$, $f'(1)=0$ and $f''(1)=-3$.
        Classify $x=1$ and give the tangent line there.
  \item Sketch a single curve with all of: domain $\R\setminus\{2\}$, a
        pole at $x=2$, horizontal asymptote $y=1$, a local minimum at
        $(0,0)$, and convex on $(2,\infty)$.
  \item Why can a function with a vertical asymptote at $x=a$ have a local
        maximum lying below a local minimum? Illustrate with problem 24.
\end{problems}

\begin{problems}[resume]
  \item For $C(q)=q^{3}-6q^{2}+15q$ on $q\ge0$: show $C$ is increasing
        throughout, find its inflection point, and show that marginal cost
        is minimal exactly there.
  \item For the same $C$: find the output minimising average cost and
        verify that $MC=AC$ at that output.
  \item With the same $C$ and price $p=15$, find the profit-maximising
        output and the maximal profit.
  \item For $R(q)=q(200-4q)$: discuss the curve on its economic domain and
        find the revenue-maximising quantity.
  \item For $C(q)=0.5q^{2}+2q+8$: find the minimum of $AC$ and confirm
        $MC=AC$ there.
  \item The adoption path is $y(t)=\dfrac{100}{1+\mathrm{e}^{-0.5t}}$.
        Find its asymptotes, show it is increasing, and locate its
        inflection point.
  \item Explain, using $AC'=(MC-AC)/q$, why average cost must be falling
        wherever marginal cost lies below it.
\end{problems}

\section{Solutions}

\begin{solutions}
  \item $f'=3x^{2}-12=3(x-2)(x+2)$: $+$ on $(-\infty,-2)$, $-$ on
        $(-2,2)$, $+$ on $(2,\infty)$. Local maximum $(-2,16)$, local
        minimum $(2,-16)$.
  \item $f'=4x^{3}-16x=4x(x-2)(x+2)$: $-$, $+$, $-$, $+$ on the four
        intervals cut by $-2,0,2$. Local minima $(\pm2,-16)$, local maximum
        $(0,0)$.
  \item $f'=(1-x)\mathrm{e}^{-x}$: positive for $x<1$, negative for $x>1$.
        Single local (and global) maximum $(1,1/\mathrm{e})$.
  \item $f'=\dfrac{1-x^{2}}{(x^{2}+1)^{2}}$: negative, positive, negative
        on $(-\infty,-1)$, $(-1,1)$, $(1,\infty)$. Local minimum
        $(-1,-\tfrac12)$, local maximum $(1,\tfrac12)$.
  \item $f'=1-\dfrac1x=\dfrac{x-1}{x}$: negative on $(0,1)$, positive on
        $(1,\infty)$. Local (global) minimum $(1,1)$.
  \item $f'=3x^{2}\ge0$ with equality only at $0$: the sign of $f'$ never
        changes, so $f$ increases throughout and $x=0$ is a
        \emph{saddle-like} stationary point, not an extremum. A vanishing
        derivative is necessary, never sufficient.
  \item $f'=\dfrac{x(x-2)}{(x-1)^{2}}$. The denominator is always positive,
        so the sign is that of $x(x-2)$: $+$, $-$, $-$, $+$ across
        $0,1,2$ (with $x=1$ excluded). Local maximum $(0,0)$, local minimum
        $(2,4)$; the function decreases on $(0,1)$ and on $(1,2)$.
  \item $f'=\dfrac{1}{2\sqrt{x}}-1$, zero when $\sqrt{x}=\tfrac12$, i.e.
        $x=\tfrac14$: positive before, negative after. Global maximum
        $\bigl(\tfrac14,\tfrac14\bigr)$.
  \item $f''=6x-6$: negative for $x<1$ (concave), positive for $x>1$
        (convex). Inflection $(1,-2)$.
  \item $f''=12x^{2}-8$: zero at $x=\pm\sqrt{2/3}$; convex outside,
        concave between. Inflections at
        $\bigl(\pm\sqrt{2/3},-\tfrac{20}{9}\bigr)$.
  \item $f'=-2x\mathrm{e}^{-x^{2}}$ and
        $f''=(4x^{2}-2)\mathrm{e}^{-x^{2}}$: zero at
        $x=\pm1/\sqrt2$, convex outside, concave between. Inflections at
        $\bigl(\pm1/\sqrt2,\mathrm{e}^{-1/2}\bigr)$.
  \item $f''=12x^{2}\ge0$, vanishing only at $x=0$ but \emph{not changing
        sign} there: the curve is convex on both sides, so there is no
        inflection. $f''=0$ is necessary, not sufficient.
  \item $f'=\dfrac{1-\ln x}{x^{2}}$, $f''=\dfrac{2\ln x-3}{x^{3}}$: zero at
        $\ln x=\tfrac32$, i.e. $x=\mathrm{e}^{3/2}$; concave before, convex
        after. Inflection at
        $\bigl(\mathrm{e}^{3/2},\tfrac32\mathrm{e}^{-3/2}\bigr)$.
  \item $f'=(2x-x^{2})\mathrm{e}^{-x}$,
        $f''=(x^{2}-4x+2)\mathrm{e}^{-x}$: zero at $x=2\pm\sqrt2$, convex
        outside those, concave between.
  \item Pole $x=-2$; equal degrees, so horizontal asymptote $y=1$.
  \item Poles $x=\pm2$; numerator degree smaller, so $y=0$.
  \item Pole $x=1$; degrees differ by one, and
        $f=x+1+\dfrac{1}{x-1}$, so the oblique asymptote is $y=x+1$.
  \item No real pole ($x^{2}+1>0$); degrees differ by one, and
        $x^{3}=x(x^{2}+1)-x$, so $f=x-\dfrac{x}{x^{2}+1}$ and the oblique
        asymptote is $y=x$.
  \item Vertical asymptote $x=0$ (as $x\to0^{+}$, $f\to-\infty$) and
        horizontal $y=0$ as $x\to\infty$, since $\ln x\ll x$.
  \item As $x\to+\infty$, $\mathrm{e}^{-x}\to0$, so $y=x$ is an oblique
        asymptote. As $x\to-\infty$ the exponential dominates and
        $f\to+\infty$ with no asymptote.
  \item Domain $\R$; no symmetry; $f=x^{2}(x-3)$ so roots $0$ (double,
        touching) and $3$; $f(0)=0$; ends $-\infty$ and $+\infty$. From
        problem 9's table: increasing on $(-\infty,0)$, decreasing on
        $(0,2)$, increasing on $(2,\infty)$, local maximum $(0,0)$, local
        minimum $(2,-4)$; concave then convex with inflection $(1,-2)$. See
        the figure.
  \item Domain $\R$; \emph{even}; $f=x^{2}(x^{2}-4)$ so roots $0$ (double)
        and $\pm2$; both ends $+\infty$; local maximum $(0,0)$, local
        minima $(\pm\sqrt2,-4)$; inflections
        $\bigl(\pm\sqrt{2/3},-\tfrac{20}{9}\bigr)$. A symmetric W.
  \item Domain $\R$; \emph{odd}; single root $0$; horizontal asymptote
        $y=0$ at both ends; local minimum $(-1,-\tfrac12)$, local maximum
        $(1,\tfrac12)$;
        $f''=\dfrac{2x(x^{2}-3)}{(x^{2}+1)^{3}}$ gives inflections at
        $(0,0)$ and $\bigl(\pm\sqrt3,\pm\tfrac{\sqrt3}{4}\bigr)$. See the
        figure.
  \item The worked example of the Core Rules: domain $\R\setminus\{1\}$,
        double root at $0$, pole $x=1$, oblique asymptote $y=x+1$, local
        maximum $(0,0)$, local minimum $(2,4)$, no inflection, concave left
        of the pole and convex right of it.
  \item Domain $\R$; \emph{even}; no roots ($f>0$ always); horizontal
        asymptote $y=0$ at both ends; global maximum $(0,1)$; inflections
        at $\bigl(\pm1/\sqrt2,\mathrm{e}^{-1/2}\bigr)$; convex outside
        them, concave between. The bell curve.
  \item Domain $\R$; no symmetry; single root $0$; as $x\to+\infty$,
        $f\to0^{+}$ (asymptote $y=0$), as $x\to-\infty$, $f\to-\infty$;
        global maximum $(1,1/\mathrm{e})$; inflection
        $(2,2/\mathrm{e}^{2})$, concave before, convex after. See the
        figure.
  \item Domain $(0,\infty)$; root $x=1$; vertical asymptote $x=0$ with
        $f\to-\infty$, horizontal asymptote $y=0$ as $x\to\infty$; global
        maximum $(\mathrm{e},1/\mathrm{e})$; inflection at
        $x=\mathrm{e}^{3/2}$. See the figure.
  \item Domain $\R\setminus\{-2\}$; root $x=1$; $f(0)=-\tfrac12$; pole
        $x=-2$ ($-\infty$ from the left, $+\infty$ from the right);
        horizontal asymptote $y=1$; $f'=\dfrac{3}{(x+2)^{2}}>0$ so strictly
        increasing on each branch with \emph{no} extrema;
        $f''=\dfrac{-6}{(x+2)^{3}}$, convex left of the pole, concave right
        of it.
  \item Domain $\R$; no symmetry; double root $0$; as $x\to+\infty$,
        $f\to0^{+}$; as $x\to-\infty$, $f\to+\infty$; local minimum
        $(0,0)$, local maximum $(2,4/\mathrm{e}^{2})$; inflections at
        $x=2\pm\sqrt2$.
  \item Domain $\R\setminus\{\pm2\}$; \emph{even}; no roots; poles
        $x=\pm2$; horizontal asymptote $y=0$;
        $f'=\dfrac{-2x}{(x^{2}-4)^{2}}$, so increasing for $x<0$ and
        decreasing for $x>0$, giving a local maximum
        $\bigl(0,-\tfrac14\bigr)$ --- a maximum with a negative value,
        since the middle branch lies entirely below the axis.
  \item Domain $[0,\infty)$; roots $0$ and $1$; $f\to-\infty$ as
        $x\to\infty$; global maximum $\bigl(\tfrac14,\tfrac14\bigr)$;
        $f''=-\tfrac14x^{-3/2}<0$, so concave throughout with no
        inflection; the tangent at $0$ is vertical.
  \item Domain $\R$; $f=(x-1)^{2}(x+2)$, so a double root at $1$ (touching)
        and a simple root at $-2$; $f(0)=2$; ends $\mp\infty$; local
        maximum $(-1,4)$, local minimum $(1,0)$; inflection $(0,2)$.
  \item It rises, but ever more slowly: increasing and concave, bending
        away from its own tangent downwards --- the shape of a production
        function with diminishing returns.
  \item $f'<0$ between the roots and $>0$ outside them, so $f$ increases on
        $(-\infty,-1)$, decreases on $(-1,3)$, increases on $(3,\infty)$: a
        local maximum at $x=-1$ and a local minimum at $x=3$.
  \item $f'(x)=kx$ with $k>0$ gives $f(x)=\tfrac{k}{2}x^{2}+c$: an upward
        parabola with vertex at $x=0$, convex everywhere.
  \item $f'(1)=0$ with $f''(1)<0$: a local \textbf{maximum} at $(1,2)$ by
        the second-derivative test. The tangent is horizontal, $y=2$.
  \item One valid answer: $f(x)=1-\dfrac{4}{(x-2)^{2}}$ has the pole and the
        asymptote $y=1$ but a maximum, not a minimum, so adjust ---
        $f(x)=\dfrac{x^{2}}{(x-2)^{2}}$ works: $f(0)=0$ is a local minimum,
        the pole is at $x=2$, and $f\to1$ at both ends. Any sketch with a
        left branch dipping to $(0,0)$ and rising to $+\infty$ at $x=2^{-}$,
        and a right branch descending convexly from $+\infty$ towards
        $y=1$, is correct.
  \item Because a local extremum is only a statement about a
        \emph{neighbourhood}. The pole cuts the domain into two intervals
        with no path between them, so the values on one side need not
        compare with those on the other. In problem 24 the maximum $(0,0)$
        lies on the left branch, the minimum $(2,4)$ on the right, and
        between them the function runs off to $\mp\infty$.
  \item $C'=3q^{2}-12q+15=3(q^{2}-4q+5)=3\bigl((q-2)^{2}+1\bigr)\ge3>0$, so
        $C$ increases for every $q$. $C''=6q-12$ vanishes at $q=2$ and
        changes sign there: inflection at $(2,14)$. Since $MC=C'$ and
        $MC'=C''$, marginal cost has its minimum exactly at $q=2$, where
        $MC=3$.
  \item $AC=q^{2}-6q+15=(q-3)^{2}+6$, minimal at $q=3$ with $AC=6$. There
        $MC=3\cdot9-12\cdot3+15=6=AC$, as the identity
        $AC'=(MC-AC)/q$ requires.
  \item $\pi=15q-(q^{3}-6q^{2}+15q)=-q^{3}+6q^{2}$. Then
        $\pi'=-3q^{2}+12q=-3q(q-4)$, zero at $q=4$ (and $q=0$), and
        $\pi''=-6q+12=-12<0$ at $q=4$: a maximum, with
        $\pi(4)=-64+96=32$.
  \item $R=200q-4q^{2}=-4(q-25)^{2}+2500$ on the economic domain
        $[0,50]$ (where the price $200-4q$ stays non-negative): a downward
        parabola, increasing on $[0,25)$, decreasing on $(25,50]$, concave
        throughout, with maximum revenue $2500$ at $q=25$ and zeros at both
        ends.
  \item $AC=0.5q+2+\dfrac{8}{q}$, so $AC'=0.5-\dfrac{8}{q^{2}}=0$ gives
        $q^{2}=16$, $q=4$, with $AC(4)=2+2+2=6$. And $MC=q+2=6$ at $q=4$:
        equal, as required.
  \item As $t\to\infty$, $y\to100$; as $t\to-\infty$, $y\to0$: two
        horizontal asymptotes. $y'=\dfrac{50\mathrm{e}^{-0.5t}}
        {(1+\mathrm{e}^{-0.5t})^{2}}>0$ always, so $y$ increases
        throughout. The inflection is where the numerator's growth turns,
        at $\mathrm{e}^{-0.5t}=1$, i.e. $t=0$, with $y=50$ --- exactly half
        the ceiling, which is the signature of the logistic curve.
  \item $q>0$, so the sign of $AC'$ is the sign of $MC-AC$. If
        $MC<AC$ then $AC'<0$ and average cost falls: producing one more
        unit costs less than the current average, which pulls the average
        down. Once $MC$ rises above $AC$ the average is pulled up, and the
        crossing point is therefore the minimum.
\end{solutions}

\begin{gallery}
\begin{tikzpicture}
  \begin{axis}[fnaxis, xmin=-1.6, xmax=3.9, ymin=-5.4, ymax=4.4,
      xtick={-1,1,2,3}, ytick={-4,-2,2}, title={$x^{3}-3x^{2}$}]
    \addplot[domain=-1.2:3.6, samples=201]{x^3-3*x^2};
    \addplot[only marks, mark=*, mark size=1.2pt, black, forget plot]
      coordinates {(0,0) (2,-4) (1,-2)};
  \end{axis}
\end{tikzpicture}\hfill
\begin{tikzpicture}
  \begin{axis}[fnaxis, xmin=-5.4, xmax=5.4, ymin=-0.75, ymax=0.75,
      xtick={-4,-2,2,4}, ytick={-0.5,0.5}, title={$x/(x^{2}+1)$}]
    \addplot[domain=-5.2:5.2, samples=201]{x/(x^2+1)};
    \addplot[only marks, mark=*, mark size=1.2pt, black, forget plot]
      coordinates {(1,0.5) (-1,-0.5) (0,0)};
  \end{axis}
\end{tikzpicture}
\end{gallery}

\begin{gallery}
\begin{tikzpicture}
  \begin{axis}[fnaxis, xmin=-1.4, xmax=5.4, ymin=-1.4, ymax=0.75,
      xtick={-1,1,2,4}, ytick={-1,0.5}, title={$x\mathrm{e}^{-x}$}]
    \addplot[domain=-1.2:5.2, samples=201]{x*exp(-x)};
    \addplot[only marks, mark=*, mark size=1.2pt, black, forget plot]
      coordinates {(1,0.3679) (2,0.2707)};
  \end{axis}
\end{tikzpicture}\hfill
\begin{tikzpicture}
  \begin{axis}[fnaxis, xmin=-0.5, xmax=9.4, ymin=-1.9, ymax=0.75,
      xtick={2,4,6,8}, ytick={-1,0.5}, title={$\ln x/x$}]
    \addplot[domain=0.14:9.2, samples=241]{ln(x)/x};
    \addplot[only marks, mark=*, mark size=1.2pt, black, forget plot]
      coordinates {(2.7183,0.3679) (4.4817,0.3347)};
  \end{axis}
\end{tikzpicture}
\gcap{Four of the completed discussions: problems 21, 23, 26 and 27, with
their extrema and inflection points marked.}
\end{gallery}

\section{Additional Problems}

\begin{problems}
  \item Monotonicity and extrema of $f(x)=2x^{3}-9x^{2}+12x$.
  \item Monotonicity and extrema of $f(x)=\dfrac{x^{2}}{x^{2}+3}$.
  \item Monotonicity and extrema of $f(x)=x\ln x$ on $x>0$.
  \item Convexity and inflections of $f(x)=x^{3}-6x^{2}+9x$.
  \item Convexity and inflections of $f(x)=\dfrac{1}{1+x^{2}}$.
  \item Convexity of $f(x)=\mathrm{e}^{x}+\mathrm{e}^{-x}$, and its
        minimum.
  \item All asymptotes of $f(x)=\dfrac{2x^{2}-3}{x+1}$.
  \item All asymptotes of $f(x)=\dfrac{x^{2}-1}{x^{2}+1}$.
  \item All asymptotes of $f(x)=x\mathrm{e}^{1/x}$.
  \item Complete discussion of $f(x)=x^{3}-6x^{2}+9x$.
  \item Complete discussion of $f(x)=\dfrac{x^{2}-4}{x^{2}}$.
  \item Complete discussion of $f(x)=(x-2)\mathrm{e}^{-x}$.
  \item Complete discussion of $f(x)=x^{2}\ln x$ on $x>0$.
  \item Complete discussion of $f(x)=\dfrac{4x}{x^{2}+4}$.
  \item Complete discussion of $f(x)=\sqrt{4-x^{2}}$.
  \item A cubic has a local maximum at $(0,4)$ and a local minimum at
        $(2,0)$. Determine it.
  \item A function has $f'(x)>0$ everywhere and $f''(x)<0$ everywhere, with
        horizontal asymptote $y=5$. Sketch it and give an example formula.
  \item Where does $f(x)=x^{4}-4x^{3}$ have inflections, and why does the
        second-derivative test fail at one of its critical points?
  \item Show that $f(x)=x+1/x$ has a local minimum at $x=1$ and a local
        maximum at $x=-1$, with the maximum value below the minimum value,
        and explain.
  \item For which $a$ does $f(x)=x^{3}+ax$ have two local extrema?
  \item For $C(q)=q^{3}-9q^{2}+30q$: find the inflection, the minimum of
        average cost, and verify $MC=AC$ there.
  \item For $C(q)=2q^{2}+50$ and price $p=40$: find the profit-maximising
        output, the profit, and the break-even outputs.
  \item A monopolist faces $p(q)=100-2q$ with $C(q)=q^{2}+20q$. Discuss the
        profit curve and find the optimum.
  \item Show that if $AC$ has a minimum at $q^{*}$ then the $AC$ curve is
        tangent to no ray from the origin except the one through
        $(q^{*},AC(q^{*}))$, and relate this to $MC=AC$.
  \item The production function $Q(L)=12L^{2}-L^{3}$ describes output from
        labour $L\in[0,12]$. Find where average product is maximal and show
        that marginal product equals average product there.
\end{problems}

\section{Additional Solutions}

\begin{solutions}
  \item $f'=6x^{2}-18x+12=6(x-1)(x-2)$: $+,-,+$. Local maximum $(1,5)$,
        local minimum $(2,4)$.
  \item $f'=\dfrac{6x}{(x^{2}+3)^{2}}$: negative for $x<0$, positive for
        $x>0$. Global minimum $(0,0)$; $f\to1$ at both ends, so no
        maximum.
  \item $f'=\ln x+1$, zero at $x=\mathrm{e}^{-1}$: negative before,
        positive after. Global minimum
        $\bigl(\mathrm{e}^{-1},-\mathrm{e}^{-1}\bigr)$.
  \item $f''=6x-12$: concave for $x<2$, convex for $x>2$; inflection
        $(2,2)$.
  \item $f'=\dfrac{-2x}{(1+x^{2})^{2}}$,
        $f''=\dfrac{6x^{2}-2}{(1+x^{2})^{3}}$: zero at
        $x=\pm1/\sqrt3$, convex outside, concave between. Inflections at
        $\bigl(\pm1/\sqrt3,\tfrac34\bigr)$.
  \item $f''=\mathrm{e}^{x}+\mathrm{e}^{-x}>0$ everywhere: convex, no
        inflection. $f'=\mathrm{e}^{x}-\mathrm{e}^{-x}=0$ at $x=0$, giving
        the global minimum $(0,2)$.
  \item Pole $x=-1$; degrees differ by one, and
        $2x^{2}-3=(x+1)(2x-2)-1$, so $f=2x-2-\dfrac{1}{x+1}$: oblique
        asymptote $y=2x-2$.
  \item No poles; equal degrees, so horizontal asymptote $y=1$ at both
        ends.
  \item As $x\to\pm\infty$, $\mathrm{e}^{1/x}\to1$ and more precisely
        $x\mathrm{e}^{1/x}=x+1+\tfrac{1}{2x}+\cdots$, so $y=x+1$ is an
        oblique asymptote. As $x\to0^{+}$, $f\to+\infty$ (vertical
        asymptote $x=0$); as $x\to0^{-}$, $f\to0$.
  \item $f=x(x-3)^{2}$: roots $0$ and $3$ (double); ends $\mp\infty$;
        $f'=3x^{2}-12x+9=3(x-1)(x-3)$ gives local maximum $(1,4)$ and local
        minimum $(3,0)$; $f''=6x-12$ gives inflection $(2,2)$; no
        symmetry, no asymptotes.
  \item $f=1-\dfrac{4}{x^{2}}$: domain $\R\setminus\{0\}$, \emph{even},
        roots $\pm2$, pole $x=0$ with $f\to-\infty$ from both sides,
        horizontal asymptote $y=1$; $f'=\dfrac{8}{x^{3}}$, so decreasing
        for $x<0$ and increasing for $x>0$ with no extrema;
        $f''=-\dfrac{24}{x^{4}}<0$: concave on both branches.
  \item Domain $\R$; root $x=2$; $f\to0^{+}$ as $x\to+\infty$
        (asymptote $y=0$) and $f\to-\infty$ as $x\to-\infty$;
        $f'=(3-x)\mathrm{e}^{-x}$ gives a maximum
        $(3,\mathrm{e}^{-3})$; $f''=(x-4)\mathrm{e}^{-x}$ gives an
        inflection $(4,2\mathrm{e}^{-4})$.
  \item Domain $(0,\infty)$; root $x=1$; $f\to0$ as $x\to0^{+}$ and
        $f\to\infty$ as $x\to\infty$; $f'=2x\ln x+x=x(2\ln x+1)$, zero at
        $x=\mathrm{e}^{-1/2}$, a minimum with value
        $-\tfrac{1}{2\mathrm{e}}$; $f''=2\ln x+3$, zero at
        $x=\mathrm{e}^{-3/2}$, an inflection.
  \item Domain $\R$; \emph{odd}; root $0$; horizontal asymptote $y=0$;
        $f'=\dfrac{4(4-x^{2})}{(x^{2}+4)^{2}}$ gives local minimum
        $(-2,-1)$ and local maximum $(2,1)$; inflections at $0$ and
        $\pm2\sqrt3$.
  \item Domain $[-2,2]$; \emph{even}; roots $\pm2$; range $[0,2]$; it is
        the upper semicircle of radius $2$, so concave throughout, maximum
        $(0,2)$, and the tangents at $x=\pm2$ are vertical. No asymptotes.
  \item Write $f=ax^{3}+bx^{2}+cx+d$. From the maximum at $(0,4)$:
        $d=4$ and $f'(0)=c=0$. From $f'(2)=12a+4b=0$ and $f(2)=8a+4b+4=0$:
        $b=-3a$ and $8a-12a+4=0$, so $a=1$, $b=-3$. Hence
        $f(x)=x^{3}-3x^{2}+4$.
  \item Increasing and concave with a ceiling: it rises steeply at first
        and flattens onto $y=5$ without reaching it. Example:
        $f(x)=5-\mathrm{e}^{-x}$.
  \item $f'=4x^{3}-12x^{2}=4x^{2}(x-3)$, critical at $x=0$ and $x=3$;
        $f''=12x^{2}-24x=12x(x-2)$, zero at $x=0$ and $x=2$ and changing
        sign at both, so inflections at $(0,0)$ and $(2,-16)$. At $x=0$ the
        second-derivative test fails because $f''(0)=0$; the sign table for
        $f'$ shows $4x^{2}(x-3)<0$ on both sides of $0$, so $x=0$ is no
        extremum --- it is the inflection.
  \item $f'=1-\dfrac{1}{x^{2}}$, zero at $x=\pm1$;
        $f''=\dfrac{2}{x^{3}}$, positive at $x=1$ (minimum, value $2$) and
        negative at $x=-1$ (maximum, value $-2$). The maximum value $-2$
        lies below the minimum value $2$ because the pole at $x=0$
        separates the two branches, so neither extremum is global.
  \item $f'=3x^{2}+a$ has two real roots exactly when $a<0$; then
        $x=\pm\sqrt{-a/3}$ and the sign of $f'$ changes at both, giving a
        maximum and a minimum. For $a\ge0$ the function is monotone
        increasing.
  \item $C''=6q-18=0$ at $q=3$: inflection, and $MC=3q^{2}-18q+30$ is
        minimal there with $MC=3$. $AC=q^{2}-9q+30=(q-4.5)^{2}+9.75$,
        minimal at $q=4.5$ with $AC=9.75$; and
        $MC(4.5)=60.75-81+30=9.75=AC$.
  \item $\pi=40q-2q^{2}-50$, so $\pi'=40-4q=0$ at $q=10$ with
        $\pi''=-4<0$: maximum profit $\pi(10)=400-200-50=150$. Break-even
        needs $\pi=0$: $2q^{2}-40q+50=0$, i.e. $q^{2}-20q+25=0$, so
        $q=10\pm\sqrt{75}$, about $1.34$ and $18.66$.
  \item $\pi=(100-2q)q-(q^{2}+20q)=80q-3q^{2}$, a downward parabola with
        zeros at $q=0$ and $q=80/3$, concave throughout. $\pi'=80-6q=0$ at
        $q=\tfrac{40}{3}\approx13.33$, with
        $\pi=\tfrac{1600}{3}\approx533.3$ and price
        $p=100-\tfrac{80}{3}\approx73.3$.
  \item The slope of the ray from the origin through $(q,C(q))$ is exactly
        $C(q)/q=AC(q)$. Minimising $AC$ therefore means finding the
        \emph{flattest} such ray, and the flattest ray that still meets the
        cost curve must touch it --- be tangent to it. At a tangency the
        ray's slope equals the curve's slope, i.e. $AC=MC$. That is the
        geometric content of $AC'=(MC-AC)/q$.
  \item $AP=Q/L=12L-L^{2}$, maximal where $12-2L=0$, i.e. $L=6$, with
        $AP=36$. And $MP=Q'=24L-3L^{2}$, so $MP(6)=144-108=36=AP$. The
        same tangency argument as for average cost applies, with the ray
        from the origin now touching the production curve.
\end{solutions}
